\section*{\S~6 Missionen}
\label{sec:missionen}
\addcontentsline{toc}{section}{\S~6 Missionen}

\begin{enumerate}[label=(\arabic*)]
    \item Die Teilnahme an Missionen ist für aktive Mitglieder verpflichtend.    
    \item Tag, Uhrzeit und Regelmäßigkeit der Mission legt die Mitgliederversammlung fest.    
    \item Die Missionen finden auf einem Arma 3 Gameserver statt.    
    \item Die Leitung kann besondere Events für bspw. Halloween, Weihnachten, Jubiläen etc. ansetzen. Die Teilnahme an diesen ist freiwillig.    
    \item Missionen müssen über die Anwesenheitsliste angekündigt werden. Diese sollte unverzüglich nach der Nachbesprechung der vorhergehenden Mission erstellt werden.   
    \item Anwesenheit, oder Abwesenheit für eine Mission müssen auf der Anwesenheitsliste vermerkt werden.
    \item Die Anwesenheitsliste wird 30 Minuten vor Missionsbeginn geschlossen. Mitglieder, die bis zu diesem Zeitpunkt keinen Vermerk über ihre Anwesenheit, oder Abwesenheit gemacht haben, gelten als unentschuldigt. Wenn innerhalb von 72 Stunden nach Missionsbeginn der Leitung erklärt wird, dass aufgrund eines Notfalls o.Ä. die Teilnahme nicht möglich war, gilt dies als entschuldigt.   
    \item Unmittelbar nach einer Mission ist eine Nachbesprechung dieser durchzuführen. In dieser soll der Missionsbauer und das eingesetzte Führungspersonal nicht nur vorstellen, was sie warum getan haben, sondern ihr Verhalten auch allem voran selbstkritisch reflektieren. Der Stab kann hierbei ebenfalls Kritik äußern. Mitglieder müssen die Möglichkeit haben sich äußern zu können. Ziel ist es Probleme frühzeitig zu erkennen, damit diese ausgebessert werden können aber auch gut funktionierende Ideen und Ansätze zu fördern, Genaueres kann die Leitung festlegen.    
    \begin{enumerate}[label=(\Alph*)]
    	\item Im Rahmen einer Nachbesprechung können, Wahlen, Abstimmungen, Beförderungen, Mitgliederversammlungen etc. satzungsgemäß durchgeführt werden.
    \end{enumerate}
    \item Die Leitung hat das Recht, einen aktiven Unteroffizier zu jedem Zeitpunkt einer Mission als Führungspersonal in dieser Mission einzusetzen, sollte es sich zeigen, dass dies notwendig ist.
    \item Die Leitung ist verpflichtet an mindestens sechs Missionstermine im Jahr Trainingsmissionen durchzuführen. Diese sollten möglichst gleichmäßig über das Jahr verteilt sein. 
       \begin{enumerate}[label=(\Alph*)]
        \item Zusätzlich zu Maßnahmen gem\"{a}{\ss} \S~4a Absatz \ref{itm:versetzung_reserve} und \hyperref[sec:umgang_mit_regelverstoeßen]{\S~8} kann die Leitung weitere disziplinarische Maßnahmen verh\"{a}ngen, wenn ein Mitglied regelm\"{a}{\ss}ig nicht an Trainings teilnimmt.
    \end{enumerate}    
\end{enumerate}
